\documentclass[11pt,reqno]{amsart}

\usepackage{amsmath,amssymb,mathtools,mathrsfs}
\usepackage{booktabs,array,tabularx,longtable}
\usepackage{enumitem,microtype}
\usepackage[hidelinks]{hyperref}
\usepackage{xurl}
\urlstyle{tt}
\allowdisplaybreaks[2]
\raggedbottom
\emergencystretch=3em
\setlist{nosep}

\newtheorem{theorem}{Theorem}[section]
\newtheorem{lemma}[theorem]{Lemma}
\newtheorem{proposition}[theorem]{Proposition}
\newtheorem{corollary}[theorem]{Corollary}
\theoremstyle{definition}
\newtheorem{definition}[theorem]{Definition}
\newtheorem{remark}[theorem]{Remark}

\newcommand{\K}{\mathbb K}
\newcommand{\Q}{\mathbb Q}
\newcommand{\F}{\mathbb F}
\newcommand{\Sym}{\operatorname{Sym}}
\newcommand{\rank}{\operatorname{rank}}
\newcommand{\im}{\operatorname{im}}
\newcommand{\Span}{\operatorname{span}}
\newcommand{\ChowRank}{\operatorname{ChowRank}}
\newcommand{\perm}{\operatorname{perm}}
\newcommand{\Gr}{\operatorname{Gr}}
\newcommand{\cD}{\mathcal D}

\title[A Certificate-Driven Proof for the Five-by-Five Permanent]
{A Certificate-Driven Proof of the Ordinary Chow Rank of the
Five-by-Five Permanent}
\author{Rongyu Xu}
\hypersetup{
  pdftitle={A Certificate-Driven Proof of the Ordinary Chow Rank of the Five-by-Five Permanent},
  pdfauthor={Rongyu Xu}}
\date{September 1, 2026}
\subjclass[2020]{Primary 15A69; Secondary 14N05, 05D05, 68V05}
\keywords{permanent, Chow rank, product rank, Koszul flattening,
computer-assisted proof, exact certificate}

\begin{document}

\begin{abstract}
Let \(\perm_5\) denote the permanent of a generic \(5\) by \(5\) matrix.
Over every field of characteristic zero, we prove that its ordinary Chow rank
is \(16\). This certificate-driven presentation separates the written
geometric argument from four finite computational lemmas. Their inputs,
exhaustive domains, arithmetic rings, expected outputs, and failure conditions
are stated in the article. The largest endpoint exhausts \(886{,}464\)
coordinate flags over each of \(\F_3,\F_5,\F_7\) and checks every extremum over
\(\Q\); independent programs also audit Glynn's identity and all
\(\binom{25}{5}\) coordinate five-planes. A deterministic fail-closed replay
links these certificates to the written proof. The theorem is
author-reviewed and computer-assisted; named external peer review and
proof-assistant formalization are not claimed.
\end{abstract}

\maketitle

\section{Introduction}

\subsection*{Notation and conventions}

Throughout the paper, \(\K\) denotes a field of characteristic zero. For a
positive integer \(n\), write \([n]=\{0,1,\ldots,n-1\}\), let \(S_n\) be the
group of bijections of \([n]\), and let \(\binom{\Omega}{a}\) denote the set
of all \(a\)-element subsets of a finite set \(\Omega\). If \(W\) is a
finite-dimensional \(\K\)-vector space, then \(W^*\), \(\Sym^aW\),
\(\Lambda^aW\), \(\Gr(a,W)\), and \(\mathbb P(W)\) denote its dual, its
degree-\(a\) symmetric and exterior powers, the Grassmannian of
\(a\)-dimensional subspaces, and the space of one-dimensional quotients,
respectively. For \(L\subset W\), the annihilator is
\(L^\perp=\{\lambda\in W^*: \lambda|_L=0\}\). We use \(\Span\), \(\im\),
\(\ker\), \(\rank\), \(\dim\), and \(\operatorname{codim}\) with their usual
linear-algebraic meanings. For a prime \(p\), \(\F_p\) is the field with
\(p\) elements.

Let
\[
 V=\Span_{\K}\{x_{ij}:(i,j)\in[5]\times[5]\},
 \qquad X=(x_{ij})_{i,j\in[5]}.
\]
For subsets \(I,J\subset[5]\) of the same cardinality, \(X_{I,J}\) is the
submatrix of \(X\) with row set \(I\) and column set \(J\), and
\(\perm(X_{I,J})\) is its permanent, defined by the same permutation formula
as below. Define
\[
 \perm_5=\sum_{\sigma\in S_5}\prod_{i=0}^4x_{i,\sigma(i)}\in\Sym^5V.
\]

For \(f\in\Sym^mV\), \(0\le a\le m\), and
\(\lambda_1,\ldots,\lambda_a\in V^*\), let
\(\partial_{\lambda_1}\cdots\partial_{\lambda_a}f\) denote iterated
directional differentiation. The order-\(a\) catalecticant and the remaining
degree-\(j\) derivative space, where \(j=m-a\), are
\[
 \begin{aligned}
 C_{a,j}(f)&:\Sym^aV^*\longrightarrow\Sym^jV,\\
 C_{a,j}(f)(\lambda_1\cdots\lambda_a)
   &:=\partial_{\lambda_1}\cdots\partial_{\lambda_a}f,\\
 \cD_j(f)&:=\im C_{a,j}(f).
 \end{aligned}
\]
For a subspace \(G\subset\Sym^jV\), we write
\(\partial G=\Span\{\partial_\lambda g:\lambda\in V^*,\ g\in G\}\).
If \(q\in\Sym^2V\), then \(M(q)\) is the symmetric matrix of the associated
bilinear form in the ordered basis \((x_{ij})\); the rank of \(q\) means
\(\rank M(q)\).

A quintic Chow term is a nonzero product of five members of \(V\), regarded
as linear forms on \(V^*\). The ordinary Chow rank
\(\ChowRank_{\K}(f)\) of \(f\in\Sym^5V\) is the least number of such terms
whose sum is \(f\); we set \(\ChowRank_{\K}(0)=0\). Within this paper we
abbreviate it to \(\ChowRank(f)\).
This is
also called product rank in the literature; see \cite{IltenTeitler2016}.

\begin{theorem}[Main theorem]\label{thm:main}
Over the field \(\K\),
\[
 \boxed{\ChowRank_{\K}(\perm_5)=16.}
\]
\end{theorem}

The proof is computer-assisted in a deliberately narrow sense. The geometric
reductions, rank inequalities, degeneration arguments, and the passage from a
hypothetical decomposition to a finite state space are written mathematical
proofs. Exact programs verify finite endpoint tables obtained from those
reductions. No random experiment, floating-point computation, SAT certificate,
Groebner-basis computation, or unpublished data is used.

The lower-bound argument has four stages. First, a Koszul flattening has rank
\(2400\) on \(\perm_5\), whereas one Chow term contributes at most \(240\).
Second, a low-rank classification gives a one-dimensional bound for the
intersection of a Chow derivative space with the permanent quadratic
derivative space. A projective incidence argument then reduces the unique
one-intersection endpoint to coordinate flags. Third, six terms are fixed in a
hypothetical decomposition of length at most fifteen. Petersen-product shadows
and a double-quotient rank inequality leave exactly \(58\) integer states.
Fourth, exact structural and finite arguments exclude all \(58\) states; the
terminal numerical contradiction is \(2215>9\cdot245=2205\).

Section~\ref{sec:foundations} develops the invariant linear algebra.
Section~\ref{sec:flag} proves the one-intersection theorem. Sections
\ref{sec:fixedsix} and \ref{sec:endpoints} carry out the fixed-six reduction
and its terminal exclusions. Section~\ref{sec:upper} gives the upper bound,
and Section~\ref{sec:replay} states the exact replay boundary.

\section{Derivative spaces and Koszul flattenings}\label{sec:foundations}

Write \(V=A\otimes B\), where \(A=\Span_{\K}\{a_0,\ldots,a_4\}\),
\(B=\Span_{\K}\{b_0,\ldots,b_4\}\), and \(x_{ij}=a_i\otimes b_j\). Put
\[
 E=\cD_2(\perm_5)=\im C_{3,2}(\perm_5),\qquad
 D=\cD_3(\perm_5)=\im C_{2,3}(\perm_5).
\]
Their standard bases are
\begin{align}
 E&=\Span\{x_{ia}x_{jb}+x_{ib}x_{ja}:i<j,\ a<b\},\label{eq:E}\\
 D&=\Span\{\perm(X_{I,J}):|I|=|J|=3\}.
\end{align}
Distinct basis elements have distinct monomial supports, so
\[
 \dim E=\binom52^2=100,\qquad \dim D=\binom53^2=100.
\]

For \(F\subset\Sym^2V\), define the first prolongation
\[
 F^{(1)}=\{g\in\Sym^3V:\partial_\lambda g\in F
 \text{ for all }\lambda\in V^*\}.
\]
The Koszul differential
\[
 \delta_F:V\otimes F\longrightarrow\Lambda^2V\otimes V,
 \qquad w\otimes uv\longmapsto
 u\wedge w\otimes v+v\wedge w\otimes u
\]
has kernel naturally isomorphic to \(F^{(1)}\). On a fixed set of three rows
and three columns, transpositions connect the six perfect matchings.
For \(f\in\Sym^5V\), define its quadratic Koszul flattening to be
\(K(f)=\delta_{\cD_2(f)}\). Consequently \(E^{(1)}=D\), and therefore
\begin{equation}\label{eq:Kperm}
 \rank K(\perm_5)=25\dim E-\dim E^{(1)}=2400.
\end{equation}

If \(T=\ell_1\cdots\ell_5\) is a Chow term, its quadratic derivative space
\(F_T:=\cD_2(T)\) is spanned by the ten two-factor products, and its first prolongation
contains the ten three-factor products. Upper semicontinuity under degeneration
gives
\begin{equation}\label{eq:single}
 \dim F_T\le10,\qquad \rank K(T)\le25\cdot10-10=240.
\end{equation}

\subsection{Quadrics of rank at most five}

\begin{lemma}[Low-rank classification]\label{lem:lowrank}
If \(0\ne q\in E\) and \(\rank M(q)\le5\), then
\[
 M(q)=U\otimes W,\qquad \rank U=\rank W=2.
\]
Here \(U\) and \(W\) are symmetric \(5\) by \(5\) matrices, and
\(\otimes\) denotes their Kronecker product.
In particular, \(\rank M(q)=4\).
\end{lemma}

\begin{proof}
Arrange \(M(q)\) as a \(5\) by \(5\) block matrix \((D_{ij})\) by rows of
variables. Each \(D_{ij}\) is symmetric with zero diagonal, and \(D_{ii}=0\).
Choose a nonzero block \(B_0\). The principal block
\(\left[\begin{smallmatrix}0&B_0\\B_0&0\end{smallmatrix}\right]\) has rank
\(2\rank B_0\). A nonzero symmetric zero-diagonal matrix cannot have rank one,
so \(\rank B_0=2\). Any connecting block with nonzero image in the cokernel of
this principal block would increase the total rank by at least two, contrary
to the rank-five hypothesis. Hence every connecting block has image in
\(\im B_0\).

A rank-two symmetric zero-diagonal matrix is
\(ab^{\mathsf T}+ba^{\mathsf T}\), with disjoint supports for \(a\) and
\(b\). Inside the fixed image \(\Span\{a,b\}\), the symmetric zero-diagonal
condition forces every block to be a scalar multiple of \(B_0\). Congruence
Schur elimination on the first two block rows leaves a symmetric zero-diagonal
complement of rank at most one, hence the complement is zero. Thus every block
is proportional to \(B_0\), so \(M(q)=U\otimes B_0\). Multiplicativity of rank and
the zero-diagonal condition give \(\rank U=\rank B_0=2\); take \(W=B_0\).
\end{proof}

\begin{corollary}[One-dimensional intersection]\label{cor:intersection}
If \(L\subset V\) and \(\dim L\le5\), then
\[
 \dim(E\cap\Sym^2L)\le1.
\]
Consequently \(\dim(E\cap F_T)\le1\) for every quintic Chow term \(T\).
\end{corollary}

\begin{proof}
If the intersection contained a projective line, Lemma~\ref{lem:lowrank}
would place that line in a Segre variety of rank-two tensor products. A line
contained in a Segre variety fixes one tensor factor. Two nonproportional
choices for the other factor have images spanning dimension at least three;
the essential-variable space then has dimension at least \(2\cdot3=6\), a
contradiction.
\end{proof}

\section{The one-intersection flag theorem}\label{sec:flag}

Let \(Q=\Sym^2V/E\), and let \(\pi:\Sym^2V\twoheadrightarrow Q\) be the
quotient map. For a subspace \(W\subset Q\), define its relative
prolongation excess by
\[
 p(W)=\dim\bigl(\pi^{-1}(W)^{(1)}\bigr)-\dim D.
\]
Thus no choice of a lift of \(W\) is involved.

\begin{theorem}[One-intersection flag theorem]\label{thm:flag}
Let \(T=\ell_1\cdots\ell_5\) be a Chow term satisfying
\[
 \dim F_T=10,\qquad \dim(E\cap F_T)=1.
\]
If \(Z:=\pi(F_T)\) is the nine-dimensional image of \(F_T\) in \(Q\), and
\[
 Z\subset W\subset Q,\qquad \dim W=10,
\]
then \(p(W)\le26\).
\end{theorem}

\begin{proof}
Extend scalars to an algebraic closure. Put
\(L_T=\Span\{\ell_1,\ldots,\ell_5\}\), with
\(d=\dim L_T\in\{4,5\}\). Consider the incidence variety in
\[
 \Gr(d,V)\times\Gr(10,\Sym^2V)\times\Gr(9,Q)
 \times\Gr(10,Q)\times\mathbb P(E)
\]
whose points \((L,F,Z',W',[q])\) satisfy
\begin{equation}\label{eq:incidence}
 F\subset\Sym^2L,\quad q\in F,\quad
 F\subset\pi^{-1}(Z'),\quad Z'\subset W'.
\end{equation}
These are zero loci of maps of universal bundles; the incidence variety is
therefore projective and stable under the row-column diagonal torus. The orbit
closure of the point defined by \(T\) contains a torus-fixed point by the Borel
fixed-point theorem \cite[Chapter~I, Theorem~4.7]{Borel1991}.
Corollary~\ref{cor:intersection} preserves the
one-dimensional intersection at the limit, while upper semicontinuity of
kernel dimension shows that it suffices to bound the relative prolongation at
the fixed point.

The quotient \(Q\) is multiplicity-free for the row-column diagonal torus.
Indeed, its one-dimensional weight spaces consist of the (25) square
weights, the (50) same-row weights, the (50) same-column weights, and one
quotient line for each of the (100) rectangles: in a rectangle the two
matching monomials have the same weight and their sum spans the corresponding
line of (E). Thus
\[
 \dim Q=25+50+50+100=225,
\]
and all these quotient weights are distinct. Consequently, at a fixed point,
\(L,F,Z',W'\), and \([q]\) are coordinate objects. The rank-four form \(q\)
is a rectangle permanent. The finite endpoint is therefore a completely
specified enumeration of coordinate four- and five-plane flags.
For each flag the program constructs the divided-power integer matrix of the
relative prolongation map. It checks
\begin{align*}
 21{,}600&\text{ four-plane flags},\\
 432{,}432&\text{ attached five-plane flags},\\
 432{,}432&\text{ external five-plane flags}.
\end{align*}
The maxima are respectively \(22,26,22\); exactly twenty attached flags attain
\(26\). The entire enumeration is repeated over \(\F_3,\F_5,\F_7\), with
identical histograms, and all twenty extrema are row-reduced over \(\Q\).
For an integer matrix \(M\), \(\rank_{\F_p}M\le\rank_{\Q}M\), hence
\(\dim\ker_{\Q}M\le\dim\ker_{\F_p}M\). Thus the finite-field computation is
used only in the valid direction, as an upper bound on the characteristic-zero
kernel. This proves \(p(W)\le26\).
\end{proof}

\begin{remark}
The endpoint enumeration is a premise of the proof. It is independently
implemented by
\path{scripts/perm5_one_intersection_independent_multifield.py}; it neither
imports the frozen generator nor reads the frozen result.
\end{remark}

\section{The fixed-six reduction}\label{sec:fixedsix}

Equations~\eqref{eq:Kperm} and \eqref{eq:single} first give
\(\ChowRank(\perm_5)\ge10\). Suppose, for contradiction, that a minimal
decomposition has length \(r\le15\). Then \(10\le r\le15\). Keep six
original nonzero terms fixed and repeatedly replace a remaining nonzero term
\(T\) by \(\frac12T+\frac12T\). This produces
\[
 \perm_5=T_1+\cdots+T_6+T_7+\cdots+T_{15},
\]
where the residual \(R=T_7+\cdots+T_{15}\) has exactly nine Chow terms. No
separate lower bound of fifteen is used.

Let \(C=T_1+\cdots+T_6\), and put
\[
 J=\im C_{2,3}(C),\quad H=\im C_{3,2}(C),\quad
 S=D\cap J,\quad G=E\cap H,
\]
\[
 s=\dim S,\qquad t=\dim G,\qquad d=\dim(H/G),\qquad h=\dim H=t+d.
\]
For \(1\le i\le6\), set \(F_i:=\cD_2(T_i)\), and write
\(U:=F_1+\cdots+F_6\).
Only \(H\subset U\) is asserted; equality is reserved for the near-maximal
states proved below.

The bases of \(D\) and \(E\) are indexed by pairs of three-subsets and pairs
of two-subsets. Complementation identifies the derivative shadow with the
open neighborhood operator in the product of two Petersen graphs. Its exact
minimum-shadow table is displayed below. The terminology follows classical
shadow theory \cite{Kruskal1963,Katona1968}, although this table is a
graph-specific exact computation.
\begin{equation}\label{eq:shadow}
\begin{array}{c|rrrrrrrrrrrr}
q&0&1&2&3&4&5&6&7&8&9&10&11\\ \hline
\mu_q&0&9&15&18&18&24&27&27&30&30&30&35\\[1mm]
q&12&13&14&15&16&17&18&19&20&21&22&23\\ \hline
\mu_q&35&36&36&36&36&42&45&45&48&48&48&52.
\end{array}
\end{equation}
Here \(\mu_q\) is the minimum size of the open neighborhood of a
\(q\)-vertex subset in the product graph. The dynamic recurrence uses only
\[
 \nu(k)=(0,3,5,6,6,8,9,9,10,10,10),\qquad0\le k\le10,
\]
where \(\nu(k)\) is the corresponding minimum-neighborhood function for one
Petersen graph. The recurrence is exact
integer arithmetic and is reconstructed by the replay.

Let \(\overline H:=\pi(H)=(E+H)/E\subset Q\), and set
\(p_d:=p(\overline H)=\dim((E+H)^{(1)}/D)\); the subscript records
\(d=\dim\overline H\). Since
\(J\subset(E+H)^{(1)}\), \(D\cap J=S\), and \(\partial S\subset G\),
\begin{equation}\label{eq:plower}
 p_d\ge t+d-s\ge \mu_s+d-s.
\end{equation}
Applying Sylvester's rank inequality to the two quadratic Koszul maps, and
using \(\rank K(C)\le24h\), gives the central double-quotient inequality
\begin{equation}\label{eq:double}
 \rank K(R)\ge2400+(25d-p_d)-(25s-\mu_s).
\end{equation}
Corollary~\ref{cor:intersection}, applied after lifting \(E\cap U\) to
\(\bigoplus_iF_i\) and projecting to five blocks, gives
\begin{equation}\label{eq:tbound}
 t\le\dim(E\cap U)\le5\cdot10+1=51.
\end{equation}
Since \(\mu_{23}=52\), we have \(s\le22\). The exact global prolongation bounds
together with \eqref{eq:plower}--\eqref{eq:tbound} exclude \(s\le18\) and
\(d\le8\). Thus
\[
 19\le s\le22,\qquad d\ge9,\qquad h\le60.
\]
For fixed \(s\), the possible \((t,d)\) form an integer triangle of side
\(53-\mu_s\). Hence their number is exactly
\[
 \binom82+3\binom52=58.
\]

\begin{lemma}[Near-maximal coupling]\label{lem:coupling}
If a fixed-six state satisfies \(h\ge57\), then \(H=U\).
\end{lemma}

\begin{proof}
Put
\[
 e=60-\sum_i\dim F_i,\quad
 k=\sum_i\dim F_i-\dim U,\quad c=\dim U-h.
\]
Then \(e+k+c=60-h\le3\). If \(c>0\) and \(k=0\), differentiation of a
cubic relation into the direct sum of the quadratic spaces, followed by
Euler's identity, makes every component zero. If \(k=1\), the derivatives of
a cubic relation lie in the unique quadratic relation; equality of mixed
partials forces a pure cube, which cannot occur in a derivative space of
quadratic dimension at least nine. If \(k=2\), either the same argument
applies or all nonzero cubic components lie in \(\Sym^3L\) for one two-plane
\(L\). A squarefree cubic space contains no nonzero binary cubic, while two
remaining components force the common quadratic relation space to contain
\(\Sym^2L\), of dimension three, contradicting \(k=2\). Finally, \(k\ge3\)
is incompatible with \(e+k+c\le3\) and \(c>0\). Therefore \(c=0\).
\end{proof}

\section{Exclusion of the 58 states}\label{sec:endpoints}

The states are partitioned by mutually exclusive integer conditions. The
route count is \(1+38+1+2+1+1+5+9=58\).
\begin{table}[ht]
\centering
\caption{Complete routing of the fixed-six states.}\label{tab:routes}
\begin{tabular}{@{}lr@{}}
\toprule
Exclusion mechanism & Number of states\\ \midrule
Petersen-triangle equality count at \((s,d,t)=(19,9,45)\) &1\\
Global \(p_d\) bound and \eqref{eq:double} &38\\
Equality geometry for \(d=9\) &1\\
The bound \(p_{11}\le55\) &2\\
The bound \(p_{12}\le80\) &1\\
The sharpened bound \(p_{12}\le61\) &1\\
Flag-annihilator rank gap for \(d=11,12\) &5\\
Three-orbit exclusion for \(d=10\) &9\\ \midrule
Total &58\\ \bottomrule
\end{tabular}
\end{table}

For auditability, the complete state list and its unique route assignment are
included below. It is generated independently by
\path{scripts/perm5_fixed_six_state_table.py}.
\input{state_table.tex}

\begin{proposition}[Non-\(d=10\) states]\label{prop:nondten}
Every one of the forty-nine states outside the \(d=10\) row is impossible.
\end{proposition}

\begin{proof}
The equality cases in the exact shadow computation underlying
\eqref{eq:shadow} identify the unique \(22\mapsto48\) row-star or column-star
flag; deleting one cubic basis vector recovers the \(21\mapsto48\) case by
parent counting. These finite claims are reconstructed by the unified replay
in Section~\ref{sec:replay}. For \(S\subset D\), \(G\subset E\),
\(\partial S\subset G\), and a five-plane \(L\subset V\), set
\[
 W_L=\Span\{\partial_\lambda s:s\in S,\ \lambda\in L^\perp\}.
\]
Here a subscript records dimension, so \(S_{22}\) and \(G_{48}\) denote the
complete flag with \(\dim S_{22}=22\) and \(\dim G_{48}=48\). On this flag,
\(\operatorname{codim}_{G_{48}}W_L\le4\); on the deleted-point flag the bound is
seven. The failure locus is a projective closed incidence. After Borel
degeneration of both the flag and \(L\), the seven row-occupancy partitions
of five variables have maximal parent-functional counts
\[
 4,4,2,1,1,0,0.
\]
The independent program streams all \(\binom{25}{5}=53{,}130\) coordinate
five-planes and reproduces the four relevant global maxima \(4,4,5,7\).

For \(1\le i\le6\), let
\(L_i:=\Span\{\ell:\ell\text{ is a factor of }T_i\}\subset V\). Define
the addition map and its relation kernel by
\[
 \mu_F:\bigoplus_{i=1}^6F_i\longrightarrow U,
 \qquad (f_1,\ldots,f_6)\longmapsto\sum_{i=1}^6f_i,
 \qquad K_F:=\ker\mu_F.
\]
If \(\operatorname{pr}_i\) is the \(i\)-th coordinate projection, put
\(R_i:=\operatorname{pr}_i(K_F)\subset F_i\). For \(u\in U\), choose a lift
\((f_1,\ldots,f_6)\in\bigoplus_iF_i\) and define
\(\rho_i(u):=f_i+R_i\in F_i/R_i\). This is independent of the lift. Writing
\(\rho=(\rho_1,\ldots,\rho_6)\), the six-block relation estimate is
\[
 \rho:U\longrightarrow\bigoplus_iF_i/R_i,\qquad
 \dim\ker\rho\le5\dim K_F.
\]
The chain-map identity makes \(\rho_i|_G\) vanish on \(W_{L_i}\). The
preceding codimension bounds and the strict comparisons, including \(43>42\)
and \(48>42\), exclude the five relative \(d=11,12\) states. The Petersen
equality count, the exact no-crossing bounds, and the global \(p_{11},p_{12}\)
graph bounds exclude the remaining forty-four states. Each finite table is
reconstructed from its displayed recurrence or graph and is covered by the
unified replay described in Section~\ref{sec:replay}.
\end{proof}

The nine remaining states satisfy
\begin{equation}\label{eq:dten}
 s\in\{20,21,22\},\quad d=10,\quad
 t\in\{48,49,50\},\quad h=t+10.
\end{equation}
Lemma~\ref{lem:coupling} gives \(H=U\). At least four of the six quadratic
blocks have dimension ten. If one such block met \(E\) in dimension one,
Theorem~\ref{thm:flag} would give \(p(W)\le26\), whereas
\eqref{eq:plower} gives at least \(36\). Thus four saturated blocks may be
chosen transverse to \(E\).

Choose a twenty-plane in \(S\) and apply simultaneous row-column Borel
compression to it and to a five-plane of variables. The coordinate limit is
a pair
\[
 \mathcal S\subset\binom{[5]}3\times\binom{[5]}3,\qquad
 \mathcal L\subset[5]\times[5],
\]
satisfying the following notation. For
\(\mathcal A\subset[5]\times[5]\), define the restricted derivative shadow
\[
 \partial_{\mathcal A}\mathcal S=
 \{(I\setminus\{i\},J\setminus\{j\}):
 (I,J)\in\mathcal S,\ (i,j)\in(I\times J)\cap\mathcal A\}.
\]
Write \(\partial\mathcal S:=\partial_{[5]\times[5]}\mathcal S\) and
\(\mathcal L^c:=([5]\times[5])\setminus\mathcal L\). The coordinate pair
satisfies
\begin{equation}\label{eq:flagbounds}
 |\mathcal S|=20,\quad |\partial\mathcal S|\le50,\quad
 |\partial_{\mathcal L^c}\mathcal S|\le40,\quad |\mathcal L|=5.
\end{equation}
Simultaneous elementary compression cannot increase either shadow. The layer
defect identity reduces \eqref{eq:flagbounds} to fourteen shifted profiles.
Inverse compression and Petersen-fibre analysis divide the profiles into three
terminal orbits, conventionally numbered \(0,1,13\).

\begin{proposition}[Terminal orbit classification]\label{prop:terminal}
For a coordinate subset \(\mathcal C\subset Q\), abbreviate
\(p(\mathcal C):=p(\Span\mathcal C)\). The fourteen shifted profiles and every inverse-compression parent satisfying
\eqref{eq:flagbounds} lie in exactly one of the terminal orbits \(0,1,13\).
In orbit \(1\), the four critical quotients are three length-two local
flag quotients and one same-row full-order quotient. Orbit \(13\)
satisfies the uniform structural bound \(p(\mathcal C)\le3|\mathcal C|\), with every crossing
marginal at most five. Orbit \(0\) reduces to a Petersen even-cycle identity
and five slice formulas.
\end{proposition}

\begin{proof}
The classification is finite but formula-driven. The orbit-\(1\) endpoint is
enumerated by a five-vertex closed formula, ten nonzero M\"obius term families,
and an eight-row table indexed by the complete graph \(K_4\) on four vertices.
The three length-two quotients are excluded by their relative local rings, and
the full-order quotient by the same-row valuation calculation. Both are
reconstructed by the replay in Section~\ref{sec:replay}. The orbit-\(13\)
estimate replaces the historical ten-subset
classification. The orbit-\(0\) formulas replace the historical list of
ninety-two forest edges. Integer programs independently reconstruct the
shifted profiles, visible-shadow parents, fibre profiles, terminal gaps, and
the orbit reductions. Their output is compared byte-for-byte with the frozen
exact payloads before the result is accepted.
\end{proof}

\begin{proposition}[Terminal Koszul contradiction]\label{prop:terminalrank}
For every \(d=10\) state, the residual polynomial
\(R=\perm_5-(T_1+\cdots+T_6)\) has Koszul rank at least \(2215\). A sum of
nine quintic Chow terms has Koszul rank at most \(2205\).
\end{proposition}

\begin{proof}
In orbit \(0\), the exact tangent graph has \(4100\) vertices and kernel
dimension eight. Linear reductivity upgrades the trivial cotangent action to
the complete local ring. Boolean-Fourier character orthogonality gives the
shortening identity and the constancy criterion required by the terminal
symbol. Permutations of the three displayed slots generate the full
\(S_4\)-symmetry; the resulting prolongation dimensions give
\(\rank K(R)\ge2215\). The orbit-\(1\) and orbit-\(13\) terminals satisfy the
same or a stronger bound by Proposition~\ref{prop:terminal} and their local
symbol estimates.

For a single residual Chow term, the relevant quadratic derivative space has
dimension at most ten, and its prolongation kernel has dimension at least five.
Hence its terminal Koszul rank is at most \(25\cdot10-5=245\). Nine terms
therefore have total rank at most \(9\cdot245=2205<2215\), a contradiction.
\end{proof}

Propositions~\ref{prop:nondten} and \ref{prop:terminalrank} exclude all
\(58\) fixed-six states. Consequently no decomposition of length at most
fifteen exists, and
\begin{equation}\label{eq:lower}
 \ChowRank(\perm_5)\ge16.
\end{equation}

\section{The upper bound}\label{sec:upper}

Glynn's identity \cite{Glynn2010} specializes to
\begin{equation}\label{eq:glynn}
 \perm_5=\frac1{16}
 \sum_{\substack{\epsilon\in\{\pm1\}^5\\ \epsilon_0=1}}
 \left(\prod_{r=0}^4\epsilon_r\right)
 \prod_{c=0}^4\left(\sum_{r=0}^4\epsilon_rx_{rc}\right).
\end{equation}
The right side has sixteen Chow terms. Walsh cancellation gives coefficient
one on the \(120\) permutation monomials and zero on the other \(3005\)
row-choice monomials. Thus \(\ChowRank(\perm_5)\le16\), which together with
\eqref{eq:lower} proves Theorem~\ref{thm:main}. The supplied independent
expansion checks all \(50{,}000\) column choices over the integers. Finally,
a decomposition over a characteristic-zero field remains a decomposition over
its algebraic closure, while \eqref{eq:glynn} is defined over the prime field.
The algebraically closed lower bound and the displayed upper bound therefore
give the asserted equality over every characteristic-zero field.

\section{Computational lemmas and certificate contract}
\label{sec:certificates}

This section isolates every finite assertion on which the proof depends. A
\emph{certificate contract} means a finite input class, an exact arithmetic
domain, an asserted output, and a verifier that terminates with failure unless
the asserted output is reproduced. Coordinate variables are ordered by
\(v(i,j)=5i+j\), and subsets and pairs are ordered lexicographically. All
matrix ranks below are computed by exact row reduction; no numerical
tolerance occurs.

\begin{lemma}[One-intersection endpoint certificate]
\label{lem:certificate-one-intersection}
For the coordinate fixed points arising in Theorem~\ref{thm:flag}, the
relative prolongation maxima for four-plane, attached five-plane, and external
five-plane flags are, respectively,
\[
 22,\qquad 26,\qquad 22.
\]
The exhaustive domain has \(21{,}600+432{,}432+432{,}432=886{,}464\) flags,
and exactly twenty attached flags attain the middle maximum.
\end{lemma}

\begin{proof}[Certificate proof]
For each flag the verifier constructs the matrix of the divided-power cubic
prolongation map in the lexicographically ordered coordinate bases. It scans
the complete finite domain over each of \(\F_3,\F_5,\F_7\), compares the three
histograms, and recomputes every extremal matrix over \(\Q\). The implementation
in \path{scripts/perm5_one_intersection_independent_multifield.py} does not
import a frozen producer or read its result. Since reduction modulo a prime
cannot increase rank, the finite-field kernel dimensions bound their rational
counterparts from above. Thus the certificate proves precisely the inequality
used in Theorem~\ref{thm:flag}.
\end{proof}

\begin{lemma}[Fixed-six state certificate]
\label{lem:certificate-fixed-six}
The integer triples \((s,d,t)\) allowed by
\eqref{eq:plower}--\eqref{eq:tbound} comprise exactly fifty-eight states. The
eight mutually exclusive routes in Table~\ref{tab:routes} contain
\(1,38,1,2,1,1,5,9\) states and exhaust this set.
\end{lemma}

\begin{proof}[Certificate proof]
The verifier reconstructs the Petersen-product neighborhood recurrence from
the eleven-entry function \(\nu\), enumerates the resulting integer triangle,
and evaluates the route predicates in their displayed order. It rejects an
unrouted state, a multiply routed state, a route count different from the
listed vector, or a total different from fifty-eight. The independent entry
point is \path{scripts/perm5_fixed_six_state_table.py}; it also emits the
human-readable table included in this article. The predicates are the
finite inequalities and equality cases invoked in
Propositions~\ref{prop:nondten} and~\ref{prop:terminal}.
\end{proof}

\begin{lemma}[Parent-table and terminal certificates]
\label{lem:certificate-terminal}
The coordinate five-plane parent maxima used in
Proposition~\ref{prop:nondten} are \(4,4,5,7\). The shifted-profile,
visible-shadow, inverse-compression, and terminal-orbit computations used in
Proposition~\ref{prop:terminalrank} reproduce the three terminal orbits
\(0,1,13\) and the lower bound \(\rank K(R)\ge2215\).
\end{lemma}

\begin{proof}[Certificate proof]
The independent parent-table audit streams, rather than stores, all
\(\binom{25}{5}=53{,}130\) coordinate five-planes and recomputes the four
maxima. The active terminal producers reconstruct their finite graphs,
profiles, orbit representatives, and exact integer or rational ranks. Their
outputs are compared with the immutable payloads identified in
\path{docs/perm5_v15_theorem_evidence_index.md}. A changed source digest,
missing payload, incomplete parent set, or unequal exact table is a failure.
The independent entry point
\path{scripts/perm5_terminal_independent_audit.py} reconstructs the tangent
graph and rational Fourier matrices directly from their definitions and
checks the terminal rank gap without importing the frozen producers.
\end{proof}

\begin{lemma}[Upper-bound certificate]
\label{lem:certificate-glynn}
The right-hand side of \eqref{eq:glynn} has coefficient one on each of the
\(120\) permutation monomials and coefficient zero on the other \(3005\)
row-choice monomials.
\end{lemma}

\begin{proof}[Certificate proof]
The independent integer program
\path{scripts/perm5_glynn_upper_bound_independent.py} expands all
\(16\cdot5^5=50{,}000\) column choices and compares the complete coefficient
map with the permanent support. This is a direct audit of the displayed
identity, not a probabilistic identity test.
\end{proof}

Lemmas~\ref{lem:certificate-one-intersection}--\ref{lem:certificate-glynn}
are finite lemmas in the logical proof: their conclusions, not the existence
of output files, are used in Sections~\ref{sec:flag}--\ref{sec:upper}.

\section{Exact replay and proof boundary}\label{sec:replay}

From the public package's \path{perm5/computation} directory, execute
\begin{verbatim}
mkdir -p evidence/public_replay
out=evidence/public_replay
python -B scripts/replay_perm5_v15_full.py \
  --mode normal \
  --receipt "$out/replay_normal.json"
python -B scripts/replay_perm5_v15_full.py \
  --mode optimized \
  --receipt "$out/replay_optimized.json"
python -B scripts/perm5_fixed_six_state_table.py \
  --output-dir evidence
python -B scripts/perm5_terminal_independent_audit.py \
  --output evidence/perm5_terminal_independent_audit.json
\end{verbatim}
The replay requires Python 3 and only its standard library. Each run verifies
the immutable v14 ZIP and its ordinary-file manifest before
and after execution, runs all twenty-eight active fifth-order producers, and
then runs three independent definition-level audits. Optimized mode uses
\texttt{-O}; no correctness condition in the added audits is implemented by a
Python \texttt{assert} statement.

The three independent audits are:
\begin{enumerate}
\item the \(50{,}000\)-choice integer expansion of Glynn's identity;
\item the \(886{,}464\)-flag scan over each of \(\F_3,\F_5,\F_7\), followed by
rational checks of all twenty extrema;
\item the streamed scan of all \(\binom{25}{5}=53{,}130\) coordinate
five-planes, reproducing the parent-table maxima \(4,4,5,7\).
\end{enumerate}
The frozen source-to-output index is
\path{docs/perm5_v15_theorem_evidence_index.md}. It records SHA-256
identities for every load-bearing producer and payload. The unified verifier
fails closed on a missing file, a mismatched digest, a changed payload, or a
nonzero producer exit.

The public proof package places the frozen reviewer packet, complete producer
tree, independent verifiers, generated receipts, and source-to-output index
beside this manuscript so that no private repository access is required.

The theorem label in this article means an author-reviewed computer-assisted
theorem with deterministic exact replay. It does not mean that a named
independent mathematician has peer reviewed the complete proof, and it does
not mean that the proof has been formalized in a proof assistant. No border
Chow-rank statement and no equality for general \(n\) is claimed.

\begin{thebibliography}{9}

\bibitem{Borel1991}
A.~Borel,
\emph{Linear algebraic groups}, second ed., Graduate Texts in Mathematics,
vol.~126, Springer-Verlag, New York, 1991,
\href{https://doi.org/10.1007/978-1-4612-0941-6}
{doi:10.1007/978-1-4612-0941-6}.

\bibitem{Glynn2010}
D.~G. Glynn,
\emph{The permanent of a square matrix},
European J. Combin. \textbf{31} (2010), no.~7, 1887--1891,
\href{https://doi.org/10.1016/j.ejc.2010.01.010}
{doi:10.1016/j.ejc.2010.01.010}.

\bibitem{IltenTeitler2016}
N.~Ilten and Z.~Teitler,
\emph{Product ranks of the \(3\times3\) determinant and permanent},
Canad. Math. Bull. \textbf{59} (2016), no.~2, 311--319,
\href{https://doi.org/10.4153/CMB-2015-076-1}
{doi:10.4153/CMB-2015-076-1}.

\bibitem{Katona1968}
G.~O.~H. Katona,
\emph{A theorem of finite sets}, in: Theory of Graphs (Proc. Colloq., Tihany,
1966) (P.~Erd\H{o}s and G.~O.~H. Katona, eds.), Academic Press, New York;
Akad\'emiai Kiad\'o, Budapest, 1968, pp.~187--207.

\bibitem{Kruskal1963}
J.~B. Kruskal,
\emph{The number of simplices in a complex}, in: Mathematical Optimization
Techniques (R.~Bellman, ed.), University of California Press, Berkeley, 1963,
pp.~251--278,
\href{https://doi.org/10.1525/9780520319875-014}
{doi:10.1525/9780520319875-014}.

\end{thebibliography}

\end{document}
